\chapter{Peritoneal Dialysis Catheter Placement}

\section*{Overview}
Peritoneal dialysis catheter placement inserts a soft tube into the abdomen that allows peritoneal dialysis to filter waste using the lining of the abdominal cavity.

\section*{Alternative Names}
PD catheter insertion; Tenckhoff catheter placement; peritoneal access

\section*{Principle}
A catheter is placed through the abdominal wall into the peritoneal cavity. Dialysate fills and drains through it, and the peritoneum acts as the semipermeable membrane for solute and fluid exchange.

\section*{History}
Peritoneal dialysis was established with the Tenckhoff catheter in the 1960s and has grown into a home-based alternative to hemodialysis.

\section*{Why It Is Done}
Performed to establish access for patients who choose peritoneal dialysis, prefer home therapy, or lack good vascular access for hemodialysis.

\section*{Is This a Primary or Secondary Test or Procedure}
Primary procedure for peritoneal dialysis access.

\section*{How It Is Performed}
Performed surgically or percutaneously under anesthesia. The catheter is placed into the pelvic region of the peritoneal cavity, tunneled through the subcutaneous tissue, and exits at a convenient site.

\section*{Preparation}
Bowel preparation, bladder emptying, and prophylaxis against infection.

\section*{Contraindications and Precautions}
Extensive abdominal adhesions, hernias, and conditions limiting the peritoneal surface are relative contraindications.

\section*{Risks}
Infection (peritonitis, exit-site infection), catheter migration or obstruction, leakage, and bleeding.

\section*{Aftercare and Recovery}
The catheter is flushed regularly; dialysis typically starts after a break-in period to allow healing.

\section*{Outcome and Follow-up}
Function is judged by dialysate flow, drainage volumes, and absence of leakage or infection.

\section*{Expected Results}
Free inflow and outflow of dialysate with no infection.

\section*{Follow-up}
Ongoing exit-site care and monitoring; repositioning or replacement for dysfunction.

\section*{References}
\begin{enumerate}
  \item MedlinePlus. Peritoneal dialysis. U.S. National Library of Medicine; 2025.
  \item Current Medical Diagnosis and Treatment. McGraw-Hill; 2025.
\end{enumerate}

\clearpage
